\section{Arithmetic preliminaries at an auxiliary log scale}
\label{sec:arithmetic}

This section supplies local asymptotic formulae and upper bounds for the
coefficient and residue weights used in the amplification.  All limits in this
section are as $B\to\infty$.  Constants may depend on explicitly fixed compact
scale ranges, on the finitely many smooth norms indicated below, and later on
the fixed regularity grid and its tolerance.  They are uniform in the integer
variables and moduli in the stated ranges.  The parameter $B$ will always be
held fixed when an inner limit in $x$ is taken elsewhere in the proof.
We use the notation $e(t)=\exp(2\pi i t)$.

Take $B$ through sufficiently large positive integers and put
\[
 T=\lfloor B^{0.32}\rfloor,\qquad P_0=B^{1000},\qquad
 \mathcal P=\{p:P_0<p\le \exp(4B)\},\qquad
 R=\frac{B}{\log P_0},\qquad \ell=\log R.
\]
In particular $R\to\infty$, $\ell/\log B\to1$, and $T<P_0$.  An integer is
\emph{rough} if it has no prime divisor at most $P_0$.  Write
$\mu_{\rm Mob}$ for the M\"obius function, $\omega$ for the number of distinct
prime divisors, and $\omega_E$ for this count restricted to a set of primes $E$.

The factors $2^{-\omega}$ turn divisor sums into averages over fair splits
of prime sets.  Set
\[
 A_0(a)=(\log P_0)R^{1/2}\mu_{\rm Mob}^2(a)2^{-\omega(a)}
              \ind_{a\ {\rm rough}},\qquad
 K_0(v)=R^{1/2}2^{-\omega_{\mathcal P}(v)}.
\]
The second definition also applies to $v\in\whZ$, and to a subset
$S\subset\mathcal P$, by taking $\omega_{\mathcal P}(S)=|S|$.
For $U\subset S\subset\mathcal P$, let $a_U=\prod_{p\in U}p$, with
$a_\varnothing=1$.  The definitions give
\begin{equation}\label{eq:arith-fair-split}
 A_0(a_U)K_0(S\setminus U)
 =(\log P_0)R\,2^{-|S|}
 =B\,2^{-|S|}.
\end{equation}
A fair split selects each prime of $S$ independently with probability
$1/2$.  Thus summing the left side against a function of $U$ gives $B$
times its fair-split expectation.

The coefficient supports used below are contained in $[1,\exp(bB)]$ for some
fixed $b<4$ once $B$ is large.  Thus every prime divisor of a rough coefficient
on these supports belongs to $\mathcal P$.

\subsection{Uniform local asymptotics}

Our local estimates have two outputs.  For $A_0$, they give a smooth mean
on multiplicative windows at log scale $B$, including additive phases with
moduli and real frequencies bounded by fixed powers of $B$.  For a product $a$
formed by selecting each $p\in\mathcal P$ independently with probability
$z/p$, where $z\in\{1/4,1/2\}$, they give a smooth density for $\log a/B$
after restriction to a unit residue class, with an absolute error useful
even on intervals of width $O(1/B)$ in that coordinate.  Both outputs follow
by partial summation
from cumulative estimates with a saving of any prescribed power of $B$.
We first recall the unrestricted Selberg--Delange estimates, then exclude
the primes at most $P_0$ uniformly as that cutoff moves.
For these two values of $z$, define
\[
 g_z(n)=\mu_{\rm Mob}^2(n)z^{\omega(n)},\qquad
 g_{z,P_0}(n)=g_z(n)\ind_{n\ {\rm rough}}.
\]

\begin{lemma}[Classical estimates used in roughness removal]
\label{lem:arith-classical-sd}
For each fixed nonnegative integer $H$ and either of the above values of $z$,
there are real constants $c_j(z)$ such that, for $Y\ge3$,
\[
 \sum_{n\le Y}g_z(n)
 =Y\sum_{j=0}^{H}c_j(z)(\log Y)^{z-1-j}
       +O_H\bigl(Y(\log Y)^{z-2-H}\bigr),
\]
where
\[
 c_0(z)=\frac1{\Gamma(z)}
          \prod_p(1+z/p)(1-1/p)^z>0.
\]
For each fixed $C,D>0$, uniformly over nonprincipal Dirichlet characters
$\chi$ of modulus $d\le(\log Y)^C$, one has
\[
 \sum_{n\le Y}g_z(n)\chi(n)\ll_{C,D}Y(\log Y)^{-D}.
\]
The constants in the second estimate, and the threshold after which it
holds, need not be effective.
\end{lemma}

\begin{proof}
The first statement is the classical fixed-order Selberg--Delange expansion;
see \cite[Theorem~1]{GranvilleKoukoulopoulos2019} and
\cite[Theorem~13.2]{Koukoulopoulos2019}. The normalization and the moving
roughness cutoff required here are recorded explicitly below. In $\Re s>1$,
\[
 \sum_n\frac{g_z(n)}{n^s}=\zeta(s)^z G_z(s),\qquad
 G_z(s)=\prod_p(1+zp^{-s})(1-p^{-s})^z.
\]
Define the local powers by their power-series logarithms.  The logarithm of
each local factor of $G_z$ is $O(p^{-2\Re s})$.  Consequently $G_z$ is analytic,
nonzero, and bounded on $\Re s\ge0.9$, with all derivatives bounded on fixed
smaller compact sets.  Near $s=1$, put $u=s-1$.  The analytic factor
\[
 \frac{(u\zeta(1+u))^zG_z(1+u)}{1+u}
\]
has a convergent Taylor series and has value $G_z(1)$ at $u=0$.
Truncated Perron inversion, followed by a contour around the cut to the left
of $1$, integrates its $j$th Taylor term against
$Y e^{u\log Y}u^{j-z}$.  Hankel's reciprocal-gamma formula gives the factor
$(\log Y)^{z-j-1}/\Gamma(z-j)$.  Taylor's remainder, integrated on that
contour, is $O_H(Y(\log Y)^{z-H-2})$.  For completeness, the contour may have
right edge $1+1/\log Y$, height $\exp(c\sqrt{\log Y})$, left edge
$1-c'/\sqrt{\log Y}$, and small circular part of radius $1/\log Y$.
The classical zero-free region for $\zeta$ permits fixed sufficiently small
$c,c'>0$; the other contour pieces and the Perron truncation error are
$O_H(Y\exp(-c''\sqrt{\log Y})(\log Y)^{O_H(1)})$ and are absorbed in the
displayed remainder.  The local expansion just described also identifies
$c_0(z)=G_z(1)/\Gamma(z)$.

We spell out the character uniformity.  The twisted series is
\[
 L(s,\chi)^zG_{z,\chi}(s),\qquad
 G_{z,\chi}(s)=
  \prod_p(1+z\chi(p)p^{-s})(1-\chi(p)p^{-s})^z.
\]
The same cancellation of the linear local term bounds $G_{z,\chi}$ uniformly
in $\chi$ on $\Re s\ge0.9$.  Write $L_Y=\log Y$ and
$\Theta=\exp(c\sqrt{L_Y})$, and use the rectangle
\[
 1-c'/\sqrt{L_Y}\le\Re s\le1+1/L_Y,\qquad
 |\operatorname{Im}s|\le\Theta.
\]
The classical Dirichlet zero-free region excludes zeros in this rectangle
except possibly a real zero of a real primitive character inducing $\chi$.
Siegel's bound states that for every fixed $\varepsilon>0$ such a zero obeys
$1-\beta\gg_\varepsilon d^{-\varepsilon}$; see
\cite[Theorems~12.3 and~12.10]{Koukoulopoulos2019} for the zero-free region
and exceptional-zero bound. Choose
$\varepsilon<1/(2C)$.  Since $d\le L_Y^C$, this distance is eventually larger
than $c'/\sqrt{L_Y}$, uniformly in the allowed moduli.  The Euler factors
removed for imprimitive characters have no zeros in $\Re s>0$.
Thus an analytic logarithm of $L(s,\chi)$ exists throughout the rectangle,
agreeing with its Euler-product logarithm on its intersection with
$\Re s>1$.

There is also a uniform polynomial bound in $L_Y$ on this rectangle.
For $s=\sigma+it$ truncate the Dirichlet series at
$V=\lceil d(2+|t|)\rceil$.  The periodic character sums have modulus at most
$d$, so partial summation bounds the tail by
$O(d(1+|s|)V^{-\sigma})$.  The initial segment is at most
$O((1+\log V)V^{\max(1-\sigma,0)})$.  Since
$\log V\ll_C\sqrt{L_Y}$ and $(1-\sigma)\log V=O_C(1)$ in the rectangle,
these estimates bound $L(s,\chi)$ by a fixed power of $L_Y$.
Because $z$ is real, the modulus of the chosen analytic power is
$|L(s,\chi)^z|=|L(s,\chi)|^z$.

The coefficients have modulus at most one.  At an endpoint
$Y\in\mathbb Z+\tfrac12$, truncated Perron therefore has error
$O(Y(1+\log Y)^2\exp(-c\sqrt{L_Y}))$; this follows also by summing its usual
$\min(1,(\Theta|\log(Y/n)|)^{-1})$ error.
Shift the Perron contour to the left edge.  There is no pole or branch cut
in the twisted case.  The new vertical integral is bounded by
$Y\exp(-c'\sqrt{L_Y})$ times a fixed power of $L_Y$, and the horizontal
integrals have the additional factor $\Theta^{-1}$.  This proves an exponential
square-root-log saving and hence the claimed saving of any fixed logarithmic
power.  For arbitrary $Y\ge3$, let $Y'$ be the least element of
$\mathbb Z+\tfrac12$ with $Y'\ge Y$.  Then $0\le Y'-Y<1$ and
$d\le(\log Y)^C\le(\log Y')^C$.  Replacing $Y$ by $Y'$ changes the sum by
at most one term and any smooth main term by an amount absorbed in the
stated errors.  This gives both estimates for arbitrary real $Y$.
\end{proof}

\begin{lemma}[The moving roughness cutoff]
\label{lem:arith-rough-sd}
Fix $D>0$.  There is a fixed integer $H=H(D)$ and real coefficients
$C_{j,z}(P_0)$, $0\le j\le H$, such that uniformly for
$B^{0.89}\le\log Y\le4B$,
\[
 \sum_{n\le Y}g_{z,P_0}(n)
 =Y\sum_{j=0}^{H}C_{j,z}(P_0)(\log Y)^{z-1-j}
       +O_D(YB^{-D}).
\]
There is an absolute constant $C_0$ such that for every fixed $j$,
\[
 |C_{j,z}(P_0)|\ll_j(\log P_0)^{C_0+j},\qquad
 C_{0,z}(P_0)=c_0(z)\prod_{p\le P_0}(1+z/p)^{-1}
       \sim\frac{e^{-\gamma_{\rm E}z}}{\Gamma(z)}(\log P_0)^{-z},
\]
where $\gamma_{\rm E}$ is Euler's constant.  Uniformly for nonprincipal
$\chi$ of modulus $d\le B^{15}$ in the same size range,
\[
 \sum_{n\le Y}g_{z,P_0}(n)\chi(n)\ll_D YB^{-D}.
\]
\end{lemma}

\begin{proof}
Define the multiplicative function $h_z$ by
$h_z(p^e)=(-z)^e$ for $p\le P_0$ and $e\ge1$, and by
$h_z(p^e)=0$ for $p>P_0$ and $e\ge1$.  Its Euler factors give the exact identity
$g_{z,P_0}=g_z*h_z$.  If $\epsilon_0=1/\log P_0$, then
\[
 \sum_v\frac{|h_z(v)|}{v^{1-\epsilon_0}}
 =\prod_{p\le P_0}(1-zp^{-1+\epsilon_0})^{-1}
       \ll(\log P_0)^{C_0}.
\]
Indeed $p^{\epsilon_0}\le e$ in the product, its logarithm is bounded by
a constant times $\sum_{p\le P_0}p^{-1+\epsilon_0}$ plus a bounded sum of
square terms, and Mertens' estimate bounds that first sum by
$O(\log\log P_0)$.  Increasing the absolute constant $C_0$ if necessary gives
the displayed assertion for both values of $z$.  Since
$(\log v)^k\le k!\epsilon_0^{-k}v^{\epsilon_0}$, it also gives
\[
 \sum_v\frac{|h_z(v)|}{v}(\log v)^k
       \ll_k(\log P_0)^{C_0+k}.
\]

Put $L_Y=\log Y$.  In the convolution sum the terms $v>\sqrt Y$ contribute,
after division by $Y$, at most
\[
 \sum_{v>\sqrt Y}\frac{|h_z(v)|}{v}
       \ll(\log P_0)^{C_0}\exp\left(-\frac{L_Y}{2\log P_0}\right),
\]
because the inner untwisted or twisted sum is at most $Y/v$ in absolute
value.  For $v\le\sqrt Y$, apply Lemma~\ref{lem:arith-classical-sd} at $Y/v$.
For the untwisted sum expand each factor by Taylor's formula:
\[
 (L_Y-\log v)^{z-1-j}
 =\sum_{k=0}^{H-j}\binom{z-1-j}{k}(-\log v)^kL_Y^{z-1-j-k}
  +O_H\bigl(L_Y^{z-2-H}(\log v)^{H-j+1}\bigr).
\]
This remainder is uniform for $0\le\log v\le L_Y/2$.
After summing with $h_z(v)/v$, the total Taylor and Selberg--Delange
remainders are at most
\[
 O_H\bigl(Y L_Y^{z-2-H}(\log P_0)^{C_0+H+1}\bigr).
\]
The moment sums in the polynomial coefficients may be extended to all $v$.
To see that the resulting error is negligible, combine
$(\log v)^k\ll_k\epsilon_0^{-k}v^{\epsilon_0/2}$ with the preceding Rankin
bound: the tail of each such moment is
$O_k((\log P_0)^{C_0+k}\exp(-L_Y/(4\log P_0)))$.
Thus explicitly
\[
 C_{\nu,z}(P_0)=
 \sum_{j+k=\nu}c_j(z)\binom{z-1-j}{k}(-1)^k
       \sum_v\frac{h_z(v)}v(\log v)^k.
\]
The moment bounds prove the claimed coefficient bounds.  Since
$\log P_0=1000\log B$ and $L_Y\ge B^{0.89}$, choosing the fixed integer $H$
sufficiently large makes all these errors $O_D(YB^{-D})$.
The zeroth moment is $\prod_{p\le P_0}(1+z/p)^{-1}$.  Combining it with the
Euler product for $c_0(z)$ gives
\[
 C_{0,z}(P_0)=\frac1{\Gamma(z)}
    \prod_{p\le P_0}(1-1/p)^z
    \prod_{p>P_0}(1+z/p)(1-1/p)^z.
\]
The second product tends to one, and Mertens' product formula proves its
stated asymptotic and positivity.

For a nonprincipal character, use the same convolution with
$h_z(v)\chi(v)$ and the second estimate of
Lemma~\ref{lem:arith-classical-sd}.  On $v\le\sqrt Y$ we have
$\log(Y/v)\ge L_Y/2$ and, for all large $B$,
$d\le B^{15}\le(\log(Y/v))^{18}$.
Take an arbitrarily large fixed saving exponent in that lemma.  The factor
$\sum_v|h_z(v)|/v\ll(\log P_0)^{C_0}$ is absorbed by increasing that exponent.
The terms $v>\sqrt Y$ have already been bounded independently of $\chi$.
This proves the required uniform character estimate.
\end{proof}

We now turn these counting estimates into local densities for the
coefficient weights and probability laws for randomly selected prime products.
Choose the expansion order once, sufficiently large to use
Lemma~\ref{lem:arith-rough-sd} with $D=200$ for both values of $z$.
For $L>0$ write
\[
 P_{z,B}(L)=\sum_{j=0}^{H}C_{j,z}(P_0)L^{z-1-j},\qquad
 D_{z,B}(L)=P_{z,B}(L)+P'_{z,B}(L).
\]
Thus $D_{z,B}(\log Y)$ is exactly the derivative with respect to $Y$ of
$YP_{z,B}(\log Y)$.  Define
\[
 m_B(s)=(\log P_0)R^{1/2}D_{1/2,B}(Bs),\qquad
 Q_{z,B}=\prod_{p\in\mathcal P}(1-z/p),\qquad
 f_{z,B}(s)=B Q_{z,B}D_{z,B}(Bs).
\]
These are finite combinations of smooth powers on $s>0$.
The coefficient bounds and $\log P_0=B^{o(1)}$ imply, for every fixed
nonnegative derivative order, convergence of these functions and their
derivatives uniformly on compact subintervals of the indicated domains:
\[
 m_B(s)\longrightarrow m(s)=c_A s^{-1/2}\quad(0<s<4),\qquad
 f_{z,B}(s)\longrightarrow c_zs^{z-1}\quad(0<s<3.3),
\]
where
\[
 c_A=\frac{e^{-\gamma_{\rm E}/2}}{\Gamma(1/2)},\qquad
 c_z=\frac{e^{-\gamma_{\rm E}z}}{4^z\Gamma(z)}.
\]
Here $Q_{z,B}\sim(4R)^{-z}$ by Mertens' estimate.
To check the derivative assertion directly, the $j=0$ term has the stated
limit after normalization.  Every term with $j\ge1$, as well as the
$P'_{z,B}$ part, gains at least one power of $B^{-1}$ relative to that term
on a fixed compact interval, with only a fixed power of $\log P_0$ lost.
The same reasoning applies after any fixed number of $s$-derivatives.
Moreover the leading term dominates uniformly for $Bs\ge B^{0.89}$:
each ratio of a lower-order term to it is bounded by a fixed power of
$\log P_0$ divided by a positive power of $Bs$.
Consequently $f_{z,B}(s)\ll s^{z-1}$ for
$B^{-1/10}\le s\le3.2$, and all these densities are positive on that range
for sufficiently large $B$.

\begin{proposition}[Local coefficient and product laws]
\label{prop:arith-local-laws}
Let $w_3$ be smooth and supported in a fixed compact subinterval of
$(0,\infty)$.  Uniformly for
$0.9B\le\log X\le2.2B$, $q\le B^{15}$, $(u,q)=1$, and real $\xi$ with
$|\xi|\le B^{14}$,
\begin{equation}
 \sum_c A_0(c)w_3(c/X)e(cu/q+\xi c/X)
 =X\frac{\mu_{\rm Mob}(q)}{\varphi(q)}
       \int w_3(t)m_B(\log(Xt)/B)e(\xi t)\,dt
       +O(XB^{-50}).
 \label{eq:1}
\end{equation}
The convention $q=1$ is included.  The constant in the error is controlled
by a fixed constant times finitely many low-order smooth norms of $w_3$;
the expansion order defining $m_B$ is independent of $w_3$.

Let $\nu_z$ be the law of the product of primes of $\mathcal P$ selected
independently with probabilities $z/p$, and write $s=\log a/B$ for its log
coordinate.  For $d\le B^{15}$, every such product is a unit modulo $d$.
Uniformly for all intervals $I'\subset[B^{-1/10},3.2]$ and unit residues
$r\pmod d$,
\begin{equation}
 \Prob_{\nu_z}(s\in I',\ a\equiv r\pmod d)
 =\frac1{\varphi(d)}\int_{I'}f_{z,B}(s)\,ds+O(B^{-80}).
 \label{eq:2}
\end{equation}
Uniformly for $0\le\delta\le1$,
\begin{equation}
 \Prob_{\nu_z}(s\le\delta)\ll(\delta+1/R)^z.
 \label{eq:3}
\end{equation}
For each positive integer $n\le\exp(3.2B)$,
\begin{equation}
 \nu_{1/2}(n)\ll\frac{A_0(n)}{Bn},
 \label{eq:4}
\end{equation}
where both sides are zero unless $n$ is a squarefree rough product, apart
from the allowed empty product $n=1$.
\end{proposition}

\begin{proof}
Since $P_0>B^{15}$, all rough integers are units for every modulus in the
statement.  Character orthogonality and
Lemma~\ref{lem:arith-rough-sd} give, for each unit class,
\[
 \sum_{\substack{n\le Y\\n\equiv r\ (d)}}g_{z,P_0}(n)
 =\frac{Y}{\varphi(d)}P_{z,B}(\log Y)+O(YB^{-200}).
\]
There is no loss of a factor $\varphi(d)$ here: it is cancelled by the
normalization in character orthogonality.  Partial summation, multiplication
by $(\log P_0)R^{1/2}$ when $z=1/2$, and summation over the unit residues
against $e(ur/q)$ now prove \eqref{eq:1}.  The sum of the latter phases is
the Ramanujan sum $\mu_{\rm Mob}(q)$ because $(u,q)=1$.
For explicit error accounting, summation over residues costs at most $q$,
the normalizing factor is at most $B$ eventually, and differentiating
$w_3(t)e(\xi t)$ costs at most a constant times
$(1+|\xi|)(\|w_3\|_\infty+\|w'_3\|_1+\|w_3\|_1)$ on the fixed support.
Thus the initial $B^{-200}$ saving exceeds all these losses by much more
than the claimed $B^{-50}$.  On that support $\log(Xt)$ is in the range of
the preceding lemma once $B$ is sufficiently large.

For a squarefree product $n$ of the allowed primes, independence gives the
exact mass
\[
 \nu_z(n)=Q_{z,B}\frac{z^{\omega(n)}}n
                     \prod_{p\mid n}(1-z/p)^{-1}.
\]
On $n\le\exp(3.3B)$ the last product is $1+O(B/P_0)$, uniformly: its
logarithm is $O(\omega(n)/P_0)=O(B/P_0)$.
In this range every rough prime factor is below the upper cutoff
$\exp(4B)$.  Hence the mass without that last product is exactly
$Q_{z,B}g_{z,P_0}(n)/n$.  Partial summation of the unit-class formula between
any two endpoints $e^{Ba},e^{Bb}$, with
$B^{-1/10}\le a\le b\le3.2$, gives
\[
 Q_{z,B}\sum_{\substack{e^{Ba}<n\le e^{Bb}\\n\equiv r\ (d)}}
       \frac{g_{z,P_0}(n)}n
 =\frac1{\varphi(d)}\int_a^b BQ_{z,B}D_{z,B}(Bs)\,ds
       +O(B^{-199}).
\]
Indeed each endpoint error is $O(B^{-200})$, and integrating the error
$O(tB^{-200})$ against $dt/t^2$ over a log interval of length $O(B)$ costs
$O(B^{-199})$; also $Q_{z,B}\le1$.
The $O(B/P_0)$ relative mass correction contributes at most $O(B/P_0)$ in
total, since the uncorrected masses are bounded above by the true
probabilities.  Endpoint atoms have exponentially small mass because
$n\ge\exp(B^{0.9})$.  This proves \eqref{eq:2} for arbitrary endpoint
conventions and arbitrarily short or partial intervals.  No relative error
for a short interval is asserted or needed.

If $\delta<1/R$, the event $s\le\delta$ forces the product to be empty, so
its probability is $Q_{z,B}\ll R^{-z}$.  If $1/R\le\delta\le1$, the event
forces all primes with $\log p>\delta B$ to be absent.  Mertens' product
estimate therefore bounds its probability by
\[
 \prod_{e^{\delta B}<p\le e^{4B}}(1-z/p)\ll\delta^z.
\]
These estimates prove \eqref{eq:3}, including $\delta=0$.
Finally, on a squarefree rough $n$ in the range of \eqref{eq:4},
\[
 \frac{\nu_{1/2}(n)}{A_0(n)/(Bn)}
 =Q_{1/2,B}R^{1/2}\prod_{p\mid n}(1-1/(2p))^{-1}
 =\frac12+o(1)
\]
uniformly, and also for $n=1$.  This proves \eqref{eq:4}; outside the
squarefree rough support both masses vanish.
\end{proof}

\subsection{Upper sieve bounds with reducing local weights}

We state precisely the classical upper-sieve input, in its event-space
form.  For a fixed positive integer $k$, choose once a bound $P_k>2k$
depending only on $k$.  For a sieve bound $Z\ge2$, let
$\mathcal R_Z\subset\{p:P_k<p\le Z\}$ be the set of retained primes.
Suppose their forbidden local densities obey $0\le\theta_p\le k/p$.
Mertens' theorem then gives the upper dimension condition
\[
 \prod_{\substack{p\in\mathcal R_Z\\y<p\le z}}(1-\theta_p)^{-1}
 \le\left(\frac{\log z}{\log y}\right)^k
                      \left(1+O_k(1/\log y)\right)\qquad(2\le y\le z).
\]
Let $\mathcal V$ be the ambient mass.  For each
squarefree $d$ whose prime factors all lie in $\mathcal R_Z$, let $\mathcal V_d$
be the mass satisfying the forbidden condition at every prime dividing $d$,
and put
\[
 r_d=\mathcal V_d-\mathcal V\prod_{p\mid d}\theta_p.
\]
Extend $r_d$ by zero to all other positive integers.  The classical
fundamental lemma of the upper sieve supplies $s_k>0$, depending only on
$k$, and upper weights of absolute value at most one, supported on squarefree
$d\le D$ with prime factors in $\mathcal R_Z$, such that, when $Z\le D^{1/s_k}$,
the mass avoiding the retained forbidden conditions is at most
\[
 C_k\mathcal V\prod_{p\in\mathcal R_Z}(1-\theta_p)
                     +\sum_{d\le D}|r_d|.
\]
The constants are uniform over the forbidden sets satisfying the displayed
dimension condition.  This is the bounded-dimension upper-sieve statement of
\cite[Theorems 2.4 and 3.6]{FordSieve2023}.

\begin{lemma}[Interval, rectangle, and random-root sieves]
\label{lem:arith-weighted-sieve}
There is a sufficiently small $c_k>0$, depending only on fixed $k$, with
the following properties for $Z\ge2$.

In an interval of length $N\ge1$, at most $k$ forbidden residues at each prime
$p\le Z$ give an upper bound
\[
 C_kN\prod_{p\le Z}(1-\theta_p)+O_k(N^{0.8})
       \qquad(Z\le N^{c_k}).
\]
In a rectangle of side lengths $N_1,N_2\ge1$, suppose that at each prime the
forbidden set is a union of at most $k$ proper affine lines.  With
$N_* =\min(N_1,N_2)$, there is the corresponding bound
\[
 C_kN_1N_2\prod_{p\le Z}(1-\theta_p)
       +O_k(N_1N_2N_*^{-0.4})
       \qquad(Z\le N_*^{c_k}).
\]
An application may omit specified primes from the restrictions, provided
it also omits their factors from the products.

Both assertions also apply to reducing local weights.  At each prime,
attach a factor in $[0,1]$ to each of at most $k$ specified residue
classes modulo $p$ in the interval case, or at most $k$ proper affine
lines in the rectangle case.  The factor is applied on its class or line
and equals one off it.  Replace $1-\theta_p$ by the residue average of
the product of these factors.
\end{lemma}

\begin{proof}
At a squarefree modulus $d$, the one-dimensional forbidden intersection
has at most $k^{\omega(d)}$ residue classes, each counted with error $O(1)$.
Its remainder is therefore $O(k^{\omega(d)})$.  In the rectangle there are
at most $k^{\omega(d)}d$ residue pairs: at each prime there are at most
$kp$ pairs, and the Chinese remainder theorem multiplies these bounds.
Each residue pair has count
\[
 \frac{N_1N_2}{d^2}
          +O\left(\frac{N_1+N_2}{d}+1\right).
\]
Thus its remainder is
$O(k^{\omega(d)}(N_1+N_2+d))$.
For fixed $k$,
$\sum_{d\le D}k^{\omega(d)}\ll_k D(1+\log D)^{O_k(1)}$;
one may obtain this by bounding $k^{\omega(d)}$ by a fixed divisor function
and applying the elementary hyperbola bound to its sum.
Use the preceding upper sieve with $D=N^{1/2}$ in an interval, or
$D=N_*^{1/2}$ in a rectangle, and choose $c_k<1/(2s_k)$.
The interval remainder is $N^{1/2+o(1)}$.
The rectangle remainder is at most
\[
 (N_1+N_2)N_*^{1/2+o(1)}+N_*^{1+o(1)}
            \ll N_1N_2N_*^{-0.4}
\]
for sufficiently large $N_*$.  Enlarge the constants to cover smaller
lengths.

The preceding sieve application temporarily omitted the primes $p\le P_k$.
For each unweighted choice of forbidden conditions, restore the factors of
those primes whose restrictions the application retains.  If one such prime
forbids the whole residue space, the fully sifted count is zero.  Otherwise
its surviving fraction $1-\theta_p$ is at
least $1/p$ in an interval and at least $1/p^2$ in a rectangle.  The product
of the reciprocals of these fractions over the bounded initial set is
therefore bounded in terms of $k$.  Enlarging $C_k$ restores all their
factors in the displayed main term, without changing the remainder.  This
argument is uniform when conditions coincide and when $Z$ includes only
part of the initial set.  Primes that an application chooses to omit remain
absent from both its restrictions and its product.

For the weighted assertion, at every prime independently choose whether
to forbid each specified condition, using probability one minus its local
weight.  Make these choices independently between conditions too, even if
some coincide.  For a fixed integer point, its probability
of avoiding all the randomly chosen conditions is exactly the product of
its local weights.  Apply the unweighted bound for each choice and take
expectations.  The constants and remainders are uniform because every
choice has at most $k$ residue conditions or proper lines.  Independence
between primes makes the expectation of the local-density product the
product of the expected local densities.  The latter is precisely the residue average stated in
the lemma.  The initial-prime factors were restored before this averaging,
so no lower bound for an averaged local factor is needed.  This proves the
weighted version, including coincident conditions.
Squarefreeness conditions and all factors in $[0,1]$ above $Z$ may be
dropped when applying an upper bound.
\end{proof}

In the remaining estimates, $n\asymp X$ and analogous notation mean
membership in an interval between fixed positive constant multiples of
the indicated scale.  Restricting to additional size or positivity
conditions only decreases the upper bounds.  They are uniform when
$(\log X)/B$ lies in any fixed compact subinterval of $(0,4)$, with constants
allowed to depend on that subinterval.  In particular, they apply throughout
$0.9B\le\log X\le2.2B$, the range used later.
In applications involving $TX$, this convention still leaves all
coefficient logarithms below $4B$.

\begin{lemma}[First and second coefficient moments]
\label{lem:arith-coefficient-moments}
On these intervals,
\begin{equation}
 X^{-1}\sum_{n\asymp X}A_0(n)\ll1,\qquad
 X^{-1}\sum_{n\asymp X}A_0(n)^2\ll B^{1/4+o(1)}.
 \label{eq:5}
\end{equation}
\end{lemma}

\begin{proof}
Choose a sufficiently small fixed $c>0$ and sieve up to $Z=\exp(cB)$,
as permitted by Lemma~\ref{lem:arith-weighted-sieve} throughout the fixed
log-size range.  For a single form, local weight zero at $p\le P_0$ and
local weight $t\in\{1/2,1/4\}$ above $P_0$ give residue averages
$1-1/p$ and $1-(1-t)/p$, respectively.  Their product is
\[
 \ll (\log P_0)^{-1}
             \left(\frac{\log Z}{\log P_0}\right)^{t-1}
       \ll_c(\log P_0)^{-1}R^{t-1}.
\]
Dropping squarefreeness only increases the sum.  Put
$H_B=(\log P_0)R^{1/2}$.  Multiplication by $H_B$ for the first moment gives
one.  Multiplication by $H_B^2$ for the second moment gives
\[
 H_B^2(\log P_0)^{-1}R^{-3/4}
       =(\log P_0)R^{1/4}=B^{1/4+o(1)}.
\]
The sieve remainders remain negligible after these polynomial factors.
The omitted primes above $Z$ up to any of the fixed log-size endpoints
have bounded reciprocal sum; in particular using this fixed small $c$
does not alter any power of $R$ in the bound.
\end{proof}

Fix a sufficiently large absolute constant $C_2$ and define, for nonzero
integers $j$,
\[
 \Sigma(j)=\prod_{p\mid j}(1+C_2/p).
\]
For every fixed $r>0$ this function satisfies
$\sum_{1\le j\le U}\Sigma(j)^r\ll_r U$ for $U\ge1$.
Indeed expand $\Sigma(j)^r=\sum_{d\mid j}b_r(d)$, with $b_r$ nonnegative,
supported on squarefree integers, and
$b_r(p)=(1+C_2/p)^r-1=O_r(1/p)$.  Then
$\sum_d b_r(d)/d=\prod_p(1+b_r(p)/p)<\infty$, proving the assertion by
summing the divisor expansion.

\begin{lemma}[Two- and three-form coefficient bounds]
\label{lem:arith-form-bounds}
Uniformly for $0<|j|\le T$ and squarefree rough $b\asymp TX$,
\begin{equation}
 \sum_{\substack{c\asymp X\\b+jc\asymp TX,\ b+jc>0}}
               A_0(c)A_0(b+jc)\ll X\Sigma(j).
 \label{eq:6}
\end{equation}
On the same fixed log-size ranges,
\begin{equation}
 \sum_{\substack{b\asymp TX,\ c\asymp X\\a=b+jc\asymp TX,\ a>0}}
               A_0(b)A_0(a)A_0(c)\ll TX^2\Sigma(j).
 \label{eq:7}
\end{equation}
The implied constants do not depend on $b,j,X$, or $B$.
\end{lemma}

\begin{proof}
Use the same sieve limit $Z=\exp(cB)$ with $c$ sufficiently small.
For \eqref{eq:6}, at $p\nmid jb$ the roots of $c$ and $b+jc$ are distinct.
At a prime at most $P_0$ the local density of avoiding them is $1-2/p$;
at a larger prime the residue average of the two weights $1/2$ is $1-1/p$.
Apart from a factor $1+O(p^{-2})$, these are the squares of the single-form
factors in the preceding proof.  Bounded small primes may be omitted.
If $p\mid j$, then $p<P_0$ and $b$ is a unit at $p$.  Only the root of $c$
is forbidden.  Its loss compared with two ordinary roots is at most
$1+O(1/p)$ outside the omitted bounded set.  The product of these losses
is covered by $\Sigma(j)$ after fixing $C_2$ sufficiently large.
If $p\mid b$, then $p>P_0$ and we may drop both local conditions; their
reciprocal sum is
\[
 \sum_{p\mid b}\frac1p\le\frac{\omega(b)}{P_0}
                   \ll\frac{B}{P_0}.
\]
Thus these exceptional primes have bounded total cost.
Lemma~\ref{lem:arith-weighted-sieve} now gives density at most
$C\Sigma(j)(\log P_0)^{-2}R^{-1}$ in an interval of length comparable to
$X$.  Its normalization $H_B^2$ proves \eqref{eq:6}.

For \eqref{eq:7}, use the rectangle in $(b,c)$ of area comparable to
$TX^2$.  At $p\nmid j$, the equations
$b=0$, $c=0$, and $b+jc=0$ are three distinct lines.
Their pairwise and triple intersections have density $p^{-2}$.
Inclusion--exclusion, also with the local reducing weights, therefore
gives the product of the three single-form averages up to
$1+O(p^{-2})$.  At $p\mid j$ the lines for $a$ and $b$ coincide; since
$p<P_0$, the loss from three ordinary rough exclusions to two is at most
$1+O(1/p)$.  Again their product is covered by $\Sigma(j)$.
The rectangular sieve gives density at most
$C\Sigma(j)(\log P_0)^{-3}R^{-3/2}$, which is cancelled by $H_B^3$.
All sieve remainders are negligible because $X$ is exponential in $B$,
whereas $T$, the normalizations, and the possible singular losses are
polynomial in $B$.  Extending a support to a rectangle or interval for the
sieve is harmless: the local congruence restrictions continue to hold on
the original positive support, and the extension is used only for an upper
bound.
\end{proof}

\subsection{Regularity cutoffs and their loss}

Fix an integer $L\ge1$ and let $\Delta=1/L$.
The grid consists of $g=i/L$, $1\le i\le L$.
Fix a small $\tau>0$ and, later, a sufficiently large $C_*>0$.
The choices of $L$ and $\tau$ remain available for the additional fixed
requirements in Section~\ref{sec:moments}: first take $L$ sufficiently
large, then take $\tau$ sufficiently small.  None of the estimates of the
present section requires an upper bound for $L$ or a positive lower bound
for $\tau$.

The prefix bounds will control competing divisor representations in
Section~\ref{sec:moments}; the tail lower bounds will make the two-split
second moment in Section~\ref{sec:amplification} summable.

\begin{definition}[Regular prime sets]
\label{def:arith-regularity}
For $S\subset\mathcal P$ put
\[
 N_g(S)=\begin{cases}
    |\{p\in S:\log p\le B^g\}|,&g<1,\\
    |S|,&g=1.
 \end{cases}
\]
Let $Y_i=2^i\log P_0$, $0\le i\le i_{\max}$, where $i_{\max}$ is the first
index with $Y_i\ge4B$.  The set $S$ is regular if both of the following
conditions hold:
\[
 (g/2-\tau)\ell\le N_g(S)\le(g/2+\tau)\ell
                       \quad\hbox{for every grid point }g,
\]
and
\[
 |\{p\in S:\log p>Y_i\}|
                 \ge0.4\log(B/Y_i)-C_*
                       \quad(0\le i\le i_{\max}).
\]
An integer or profinite integer is regular when its set of dividing primes
from $\mathcal P$ is regular.  Multiplicities are not counted in these
cutoffs.
\end{definition}

Put $A=A_0\ind_{\rm reg}$ and $K=K_0\ind_{\rm reg}$, with the analogous
definition for $K(S)$.  The lower total-count bound in the definition gives
\begin{equation}
 A(a),\ K(v),\ K(S)\le B^{h_0+\tau\log2+o(1)},\qquad
             h_0=\frac{1-\log2}{2}.
 \label{eq:8}
\end{equation}
For $A$, any prime factors outside $\mathcal P$ only decrease the weight;
the additional factor $\log P_0$ is $B^{o(1)}$.
Choose $\tau$ sufficiently small that
\begin{equation}\label{eq:arith-weight-exponents}
 2(h_0+\tau\log2)<0.32,\qquad
 1-6(h_0+\tau\log2)>0.07,\qquad
 1-5(h_0+\tau\log2)>0.229.
\end{equation}
These are compatible strict conditions, since
\[
 2h_0=0.3068528\ldots,\quad
 1-6h_0=0.0794415\ldots,\quad
 1-5h_0=0.2328679\ldots.
\]
Any further decrease of $\tau$ preserves them.

\begin{lemma}[Loss from imposing regularity]
\label{lem:arith-regularity-loss}
There is an absolute $c_3>0$ such that, for each fixed choice of the grid,
tolerance, and compact log-size ranges above, there are $C>0$ and a function
$\epsilon_B\to0$ such that the part of the
untruncated sum in \eqref{eq:7} where at least one coefficient is not
regular is at most
\begin{equation}
 C TX^2\Sigma(j)\bigl(\epsilon_B+e^{-c_3C_*}\bigr).
 \label{eq:9}
\end{equation}
The constant $C$ may depend on the fixed grid, tolerance, and scale ranges,
but not on $C_*,j,X,B$.  The function $\epsilon_B$ may depend on the same
fixed data but is independent of $C_*$.  The analogous bound with scale $X$
applies to the first-moment sum of $A_0$ in \eqref{eq:5}.  For the product of
two independent coefficient
weights in a box, it applies with the corresponding area in place of
$TX^2\Sigma(j)$.

The same upper probability $O(\epsilon_B+e^{-c_3C_*})$ holds for failure of
regularity in a set of independent prime indicators on $\mathcal P$ with
parameters $1/(2p)+O(p^{-2})$, uniformly when the constant in the $O$ term
is fixed.  It also holds after omitting a deterministic subset
$E_B\subset\mathcal P$ satisfying
\[
 \sum_{p\in E_B}\frac1p\le C_E\frac{B}{P_0}
\]
for a fixed $C_E>0$.  The subset may depend on $B$, and the bound is uniform
over all such subsets for each fixed $C_E$.
\end{lemma}

\begin{proof}
We give direct upper bounds relative to the scales in the statement; no
division by the actual mass of a weighted sum is used.
Sieve with the fixed $Z=\exp(cB)$ used in
Lemmas~\ref{lem:arith-coefficient-moments} and
\ref{lem:arith-form-bounds}.  For a coefficient whose count is being tested,
replace its local weight $1/2$ on a subset $\mathcal Q\subset\mathcal P$ by
$(1/2)e^s$ or $(1/2)e^{-s}$, where $s>0$ is fixed and small enough that
$(1/2)e^s<1$.  The random-root sieve still applies.  At ordinary primes the
logarithm of its product, relative to the unmodified single-form factor,
changes by
\[
 \frac12(e^{\pm s}-1)\sum_{\substack{p\in\mathcal Q\\p\le Z}}\frac1p+O(1).
\]
The $O(1)$ is uniform for small fixed $s$.  In the three-form rectangle,
line intersections alter it by a convergent sum of $O(p^{-2})$;
the factors at primes dividing $j$ are still bounded by $\Sigma(j)$,
independently of such $s$.  The one-form case and independent two-variable
case have the same conclusion with their respective scale bounds.
All modified local weights lie in $[0,1]$, so dropping those above $Z$
preserves the direction of the bound, even for the positive tilt.

For a prefix $g<1$, take
$\mathcal Q=\{p\in\mathcal P:\log p\le B^g\}$; for $g=1$ take all of
$\mathcal P$.  Mertens' estimate gives
\[
 \sum_{\substack{p\in\mathcal Q\\p\le Z}}\frac1p=g\ell+o(\ell).
\]
For example, for $g<1$ the left side is
$g\log B-\log\log P_0+O(1)$ once $B$ is large, and its difference from
$g\ell$ is $O(\log\log B)=o(\ell)$.  There are only fixed finitely many
grid points.  The exponential Markov factor for a count exceeding
$(g/2+\tau)\ell$ is $e^{-s(g/2+\tau)\ell}$; the tilted sieve therefore
bounds its weighted contribution by the appropriate base scale times
\[
 O\left(\exp\left(
  \left[-s\tau+\frac g2(e^s-1-s)\right]\ell+o(\ell)\right)\right).
\]
For a count below $(g/2-\tau)\ell$ the corresponding bound is
\[
 O\left(\exp\left(
  \left[-s\tau+\frac g2(e^{-s}-1+s)\right]\ell+o(\ell)\right)\right).
\]
If the latter threshold is negative the event is empty.  Otherwise choose
$s$ sufficiently small in terms of the fixed $\tau$; since both exponential
remainders in brackets are $O(s^2)$ and $g\le1$, both brackets are at most
$-s\tau/2$.  Thus every prefix violation saves a fixed positive power of
$B$.  A union bound over the grid and the at most three coefficients
contributes $o(1)$ times the scale bound.

For a tail endpoint $Y=Y_i<B$, use the negative tilt on
$\mathcal Q=\{p\in\mathcal P:\log p>Y\}$.  Uniformly in these endpoints,
\[
 \sum_{\substack{p\in\mathcal Q\\p\le Z}}\frac1p
                  =\log(B/Y)+O_c(1).
\]
When $Y\le cB$ this is Mertens' estimate with upper endpoint $e^{cB}$.
When $cB<Y<B$ the sum is zero and $\log(B/Y)=O_c(1)$, which gives the same
formula.  Exponential Markov at the threshold
$0.4\log(B/Y)-C_*$ gives, relative to the appropriate base scale,
\[
 O_c\left(\exp\left(-sC_*+
       \left[0.4s+\frac12(e^{-s}-1)\right]\log(B/Y)\right)\right).
\]
Now choose a small absolute $s>0$, independently of the grid.  The bracket
is $-0.1s+O(s^2)$ and is negative.  With
$a_s=\frac12(1-e^{-s})-0.4s>0$, the sum of these bounds over the dyadic
endpoints $Y<B$ is
\[
 O_c(e^{-sC_*})\sum_{Y_i<B}(Y_i/B)^{a_s}
                         \ll_c e^{-sC_*}.
\]
For $Y\ge B$ the required lower threshold is nonpositive, so there is no
failure.  Taking $c_3=s$ proves the asserted tail loss.

The sieve remainders remain negligible in each application: their savings
are powers of intervals exponential in $B$, whereas normalization and
Markov factors are only fixed powers of $B$, and there are $O(\log B)$
tail endpoints.  For $C_*\ge0$, the Markov factors involving $C_*$ are at
most their values at zero.  Thus the residual function $\epsilon_B$ can
be chosen independently of $C_*$.  The prefix factors depend only on the
fixed grid and $\tau$.  This proves \eqref{eq:9} and its one- and two-weight
variants.

Finally let $X_p$ be the independent indicators in the probability
statement, with parameters $\lambda_p=1/(2p)+O(p^{-2})$.
For either sign and any subset $\mathcal Q$,
\[
 \E\exp\left(\pm s\sum_{p\in\mathcal Q}X_p\right)
 =\prod_{p\in\mathcal Q}(1+\lambda_p(e^{\pm s}-1)).
\]
Its logarithm is
$(e^{\pm s}-1)\sum_{p\in\mathcal Q}\lambda_p+O(\sum_p\lambda_p^2)$.
The latter error is bounded uniformly, while the parameter sums are
$g\ell/2+o(\ell)$ for a prefix and
$\tfrac12\log(B/Y)+O(1)$ for a tail with $Y<B$.
Omitting $E_B$ changes any parameter sum by at most
$O_{C_E}(B/P_0)=o(1)$, uniformly for fixed $C_E$.  Apply exactly the two
Markov calculations above, now to this product moment-generating function
and with no sieve remainder.  This
proves all the probability assertions.
\end{proof}
